% Options for packages loaded elsewhere
\PassOptionsToPackage{unicode}{hyperref}
\PassOptionsToPackage{hyphens}{url}
\PassOptionsToPackage{dvipsnames,svgnames,x11names}{xcolor}
\documentclass[
  10pt,
  fontset=fandol]{ctexart}
\usepackage{xcolor}
\usepackage[margin=16mm]{geometry}
\usepackage{amsmath,amssymb}
\setcounter{secnumdepth}{-\maxdimen} % remove section numbering
\usepackage{iftex}
\ifPDFTeX
  \usepackage[T1]{fontenc}
  \usepackage[utf8]{inputenc}
  \usepackage{textcomp} % provide euro and other symbols
\else % if luatex or xetex
  \usepackage{unicode-math} % this also loads fontspec
  \defaultfontfeatures{Scale=MatchLowercase}
  \defaultfontfeatures[\rmfamily]{Ligatures=TeX,Scale=1}
\fi
\usepackage{lmodern}
\ifPDFTeX\else
  % xetex/luatex font selection
\fi
% Use upquote if available, for straight quotes in verbatim environments
\IfFileExists{upquote.sty}{\usepackage{upquote}}{}
\IfFileExists{microtype.sty}{% use microtype if available
  \usepackage[]{microtype}
  \UseMicrotypeSet[protrusion]{basicmath} % disable protrusion for tt fonts
}{}
\makeatletter
\@ifundefined{KOMAClassName}{% if non-KOMA class
  \IfFileExists{parskip.sty}{%
    \usepackage{parskip}
  }{% else
    \setlength{\parindent}{0pt}
    \setlength{\parskip}{6pt plus 2pt minus 1pt}}
}{% if KOMA class
  \KOMAoptions{parskip=half}}
\makeatother
\setlength{\emergencystretch}{3em} % prevent overfull lines
\providecommand{\tightlist}{%
  \setlength{\itemsep}{0pt}\setlength{\parskip}{0pt}}
\usepackage{amsmath,amssymb,mathtools,needspace}
\xeCJKDeclareCharClass{CJK}{"2032,"0400->"04FF,"2160->"216B,"2460->"2473,"25A0,"25C7,"25CB,"25CF}
\IfFontExistsTF{Arial Unicode MS}{\xeCJKsetup{AutoFallBack=true}\setCJKfallbackfamilyfont{\CJKrmdefault}{Arial Unicode MS}}{}
\setcounter{secnumdepth}{0}
\setlength{\emergencystretch}{2em}
\allowdisplaybreaks
\usepackage{bookmark}
\IfFileExists{xurl.sty}{\usepackage{xurl}}{} % add URL line breaks if available
\urlstyle{same}
\hypersetup{
  pdftitle={小结},
  colorlinks=true,
  linkcolor={Maroon},
  filecolor={Maroon},
  citecolor={Blue},
  urlcolor={Blue},
  pdfcreator={LaTeX via pandoc}}

\title{小结}
\author{}
\date{}

\begin{document}
\maketitle

\subsection{小结}\label{ux5c0fux7ed3}

函数逼近是数学中一个很重要的课题，本章只介绍了函数逼近最基本的概念及常用的正交多项式，并从函数计算观点介绍了用多项式逼近连续函数的主要方法。这些都是计算数学的基础，需要详细了解的读者可参看文献 {[}4{]}。本章对函数逼近与计算的另一重要课题------有理逼近没有涉及，这方面内容可参看文献 {[}5{]}。

本章介绍的曲线拟合的最小二乘法和离散 Fourier 变换都可从最佳平方逼近得到，它们在实际中都有广泛的应用，本章着重介绍了在计算机上有效和节省时间的算法，而不要求理论上的详细讨论。特别是 FFT 算法，不但应用中意义重大，而且也是计算数学中节省运算次数的一个典范。本章介绍的改进 FFT 算法比目前使用的 FFT 算法计算次数又减少了，程序也较简单，这一算法的推导与应用可参看文献 {[}6{]}。

\begin{center}\rule{0.5\linewidth}{0.5pt}\end{center}

\subsection{相关原页校注（agent 补充）}\label{ux76f8ux5173ux539fux9875ux6821ux6ce8agent-ux8865ux5145}

以下校注来自本节覆盖的原页，可能也涉及同页相邻小节；原文照录与校正说明分开保留。

\subsubsection{原书 PDF 页 90 的校注}\label{ux539fux4e66-pdf-ux9875-90-ux7684ux6821ux6ce8}

\href{../pages/pdf-090.md}{完整原页转录} · \href{../../source-images/pdf-090.jpeg}{原页图像}

校注记录：ch03b-E020, ch03b-E021。

\begin{quote}
校注（agent 补充）：步 5、步 6 的差式中，第二个读取地址的偏移量原书均印为 \(2^{q-1}\)，与式 (3.7.16) 的 \(2^{p-1}\) 不一致；正文保留原式。步 6 的 \(j\) 循环上界原印为 \(2^{q-1}\)，相较式 (3.7.16) 与步 5 少了``\(-1\)''。按该上界执行会多计算一个 \(j\)，在 \(q=p\) 时也会越出 \(N\) 个数组单元范围。见 \href{../assets/ch03b/p090-fft-steps.png}{原程序步骤裁切}。
\end{quote}

\subsection{本节覆盖原页的校注记录}\label{ux672cux8282ux8986ux76d6ux539fux9875ux7684ux6821ux6ce8ux8bb0ux5f55}

下列记录按原页关联，含同页相邻小节的校注；计算时同时读取上文校注和完整原页。

\begin{itemize}
\tightlist
\item
  \texttt{ch03b-E020}：原书 PDF 页 90
\item
  \texttt{ch03b-E021}：原书 PDF 页 90
\end{itemize}

\end{document}
